\section{Introducción}
La ineficiencia en la gestión interna es un problema persistente,
especialmente en el entorno académico, donde la falta de formas
digitales unificadas para registrar la asistencia y el uso de servi-
cios genera información incompleta o, peor aún, incorrecta. Esta
situación es crítica, ya que la toma de decisiones administrativas
y logísticas se ve obstaculizada, y el cumplimiento de la Guía de
Registro Calificado del Ministerio de Educación, que exige datos
transversales y fiables, se hace cuesta arriba.

Por tal motivo, el diseño y la implementación de sistemas basa-
dos en códigos QR y RFID se centra en objetivos claros y cuantifi-
cables:
Reducir drásticamente el tiempo invertido en el registro de asis-
tencia.
Eliminar la incoherencia y el riesgo de pérdida de datos que aca-
rrea el registro manual en papel.
Proveer una plataforma que concentre, centralice y ofrezca aná-
lisis en tiempo real y de calidad, esencial para la gestión de recursos
humanos y físicos.

Se busca reemplazar el proceso tedioso con una herramienta que
no solo simplifique la tarea del profesor o el administrador, sino
que también genere un reporte de asistencia completo y confiable,
mejorando la seguridad y la logística interna.

\subsection{Contexto del Problema}
En las instituciones educativas colombianas, la gestión de asistencia y control de acceso representa un desafío significativo. Según datos del Ministerio de Educación Nacional, más del 60\% de las universidades aún utilizan métodos manuales o semi-manuales para el registro de asistencia, lo que genera:

\begin{itemize}
  \item Pérdida de tiempo administrativo equivalente a 8 horas semanales por docente.
  \item Errores en el registro que afectan la acreditación institucional.
  \item Dificultades en el seguimiento de indicadores de calidad educativa.
  \item Riesgos de seguridad en el campus universitario.
\end{itemize}

\subsection{Justificación Tecnológica}
La convergencia de tecnologías móviles, internet de las cosas (IoT) y computación en la nube ha creado un ecosistema propicio para la transformación digital de procesos administrativos. Los códigos QR y RFID representan soluciones maduras y accesibles que pueden ser implementadas con inversión moderada y retorno de inversión rápido.

\subsection{Objetivos del Estudio}
Este trabajo persigue los siguientes objetivos específicos:

\begin{enumerate}
  \item Diseñar e implementar un sistema integrado de control de acceso y asistencia utilizando tecnologías QR y RFID.
  \item Evaluar la efectividad de ambas tecnologías en términos de usabilidad, seguridad y rendimiento.
  \item Desarrollar una plataforma que permita la gestión centralizada de datos de asistencia.
  \item Validar la solución en un entorno universitario real con más de 5000 usuarios.
  \item Proponer un modelo de adopción tecnológica replicable en otras instituciones.
\end{enumerate}

\subsection{Estructura del Documento}
El presente documento se estructura de la siguiente manera: la Sección 2 presenta el marco teórico y trabajos relacionados; la Sección 3 describe la metodología de investigación aplicada; la Sección 4 detalla la implementación del software; la Sección 5 presenta los resultados y evaluación; la Sección 6 discute los hallazgos; finalmente, la Sección 7 presenta las conclusiones y trabajo futuro.